\chapter{Simulation Studies}

\section{Purpose of Simulation Studies}

The simulation studies will evaluate both the physical behaviour of the aircraft model and the performance of the NMPC autopilot. The studies should answer whether the model can be trimmed, whether its responses are qualitatively reasonable and whether the controller can track references without violating actuator or flight-envelope constraints.

\section{Planned Scenarios}

The planned scenarios include:
\begin{enumerate}
    \item altitude hold after an initial perturbation,
    \item commanded altitude change,
    \item commanded airspeed change,
    \item heading hold and heading change,
    \item simplified lateral path tracking,
    \item descent along a prescribed vertical path,
    \item localiser and glideslope capture,
    \item asymmetric thrust or engine-response perturbation,
    \item changed aircraft mass,
    \item actuator saturation and rate-limit stress test.
\end{enumerate}

\section{Evaluation Metrics}

The following metrics will be used to assess simulation results:
\begin{itemize}
    \item steady-state tracking error,
    \item overshoot,
    \item settling time,
    \item maximum control-surface deflection,
    \item maximum control-surface rate,
    \item maximum thrust command,
    \item constraint violation magnitude and duration,
    \item smoothness of the response,
    \item numerical robustness of the simulation.
\end{itemize}
